% Options for packages loaded elsewhere
\PassOptionsToPackage{unicode}{hyperref}
\PassOptionsToPackage{hyphens}{url}
\PassOptionsToPackage{space}{xeCJK}
\documentclass[
  chinese,
  11pt,
]{ctexart}
\usepackage{xcolor}
\usepackage[margin=2.2cm]{geometry}
\usepackage{amsmath,amssymb}
\setcounter{secnumdepth}{5}
\usepackage{iftex}
\ifPDFTeX
  \usepackage[T1]{fontenc}
  \usepackage[utf8]{inputenc}
  \usepackage{textcomp} % provide euro and other symbols
\else % if luatex or xetex
  \usepackage{unicode-math} % this also loads fontspec
  \defaultfontfeatures{Scale=MatchLowercase}
  \defaultfontfeatures[\rmfamily]{Ligatures=TeX,Scale=1}
\fi
\usepackage{lmodern}
\ifPDFTeX\else
  % xetex/luatex font selection
  \ifXeTeX
    \usepackage{xeCJK}
    \setCJKmainfont[]{Noto Serif SC}
    \setCJKsansfont[]{Noto Sans SC}
    \setCJKmonofont[]{DejaVu Sans Mono}
  \fi
  \ifLuaTeX
    \usepackage[]{luatexja-fontspec}
    \setmainjfont[]{Noto Serif SC}
  \fi
\fi
% Use upquote if available, for straight quotes in verbatim environments
\IfFileExists{upquote.sty}{\usepackage{upquote}}{}
\IfFileExists{microtype.sty}{% use microtype if available
  \usepackage[]{microtype}
  \UseMicrotypeSet[protrusion]{basicmath} % disable protrusion for tt fonts
}{}
\makeatletter
\@ifundefined{KOMAClassName}{% if non-KOMA class
  \IfFileExists{parskip.sty}{%
    \usepackage{parskip}
  }{% else
    \setlength{\parindent}{0pt}
    \setlength{\parskip}{6pt plus 2pt minus 1pt}}
}{% if KOMA class
  \KOMAoptions{parskip=half}}
\makeatother
\usepackage{longtable,booktabs,array}
\newcounter{none} % for unnumbered tables
\usepackage{calc} % for calculating minipage widths
% Correct order of tables after \paragraph or \subparagraph
\usepackage{etoolbox}
\makeatletter
\patchcmd\longtable{\par}{\if@noskipsec\mbox{}\fi\par}{}{}
\makeatother
% Allow footnotes in longtable head/foot
\IfFileExists{footnotehyper.sty}{\usepackage{footnotehyper}}{\usepackage{footnote}}
\makesavenoteenv{longtable}
\ifLuaTeX
\usepackage[bidi=basic,shorthands=off,provide=*]{babel}
\else
\usepackage[bidi=default,shorthands=off,provide=*]{babel}
\fi
\ifLuaTeX
  \usepackage{selnolig} % disable illegal ligatures
\fi
\setlength{\emergencystretch}{3em} % prevent overfull lines
\providecommand{\tightlist}{%
  \setlength{\itemsep}{0pt}\setlength{\parskip}{0pt}}
\usepackage{bookmark}
\IfFileExists{xurl.sty}{\usepackage{xurl}}{} % add URL line breaks if available
\urlstyle{same}
\hypersetup{
  pdftitle={课题10-长程Agent},
  pdflang={zh-CN},
  hidelinks,
  pdfcreator={LaTeX via pandoc}}

\title{课题10-长程Agent}
\author{}
\date{}

\begin{document}
\maketitle

{
\setcounter{tocdepth}{2}
\tableofcontents
}
\section{课题10 开题简报 · 长程 Agent 与 self-judge}\label{ux8bfeux989810-ux5f00ux9898ux7b80ux62a5-ux957fux7a0b-agent-ux4e0e-self-judge}

\begin{quote}
模块：\textbf{M8} ｜ 周次：\textbf{W14（12-14→12-20）} ｜ 算力：\textbf{T2 Kaggle 2×T4，计划 8 GPU·小时} 唐杰作业对应：\textbf{⑤ 训长程 Agent，鼓励 self-judge loop} ｜ 判分来源：\textbf{自建}（消融实验 + 轨迹日志） 手写：\textbf{R10（self-judge 协议设计）}
\end{quote}

\subsection{一、研究背景}\label{ux4e00ux7814ux7a76ux80ccux666f}

从''回答问题''到''完成任务''的跨越，靠的是 \textbf{harness（外壳）}：工具调用、上下文管理、失败重试、自我评估。 Karpathy 与李宏毅都专讲 \textbf{context engineering}；nanochat 里 \texttt{execution.py} 就是一个最小的''能执行代码的 Agent 外壳''。 本课题的落点是\textbf{长程任务}：多步串联后误差累积，以及\textbf{self-judge（自我评判）到底有没有用}。

\subsection{二、研究问题}\label{ux4e8cux7814ux7a76ux95eeux9898}

\begin{quote}
\textbf{多步任务中，失败主要来自模型能力还是来自 harness 设计？加上 self-judge loop 后成功率提升多少------还是反而下降？}
\end{quote}

\subsection{三、攻关线索}\label{ux4e09ux653bux5173ux7ebfux7d22}

\begin{enumerate}
\def\labelenumi{\arabic{enumi}.}
\tightlist
\item
  \textbf{最小 ReAct}：不用任何框架，手写 \texttt{Thought\ →\ Action\ →\ Observation} 循环（三四十行就能跑）
\item
  \textbf{工具设计}：先给两个工具（计算器 / 受限代码执行）。\textbf{工具接口的清晰度往往比模型智力更决定成败}
\item
  \textbf{长程任务构造}：把课题09 的题改造成''需要 3--5 步''的复合任务（如：先算 → 再验 → 再格式化）
\item
  \textbf{失败分类}：格式错 / 工具调错 / 目标漂移（忘了原始要求）/ 循环打转 / 提前放弃
\item
  \textbf{self-judge 的两种形态}：

  \begin{itemize}
  \tightlist
  \item
    \textbf{自评重试}：模型给自己打分，低于阈值就重做（rasbt ch5 的 self-refinement）
  \item
    \textbf{自评过滤}：生成多个候选，自己挑最好的（best-of-N 的变体）
  \end{itemize}
\item
  \textbf{关键实验（消融）}：同样的任务集，三组对照 = 无 self-judge / 自评重试 / 自评过滤
\item
  \textbf{警惕}：self-judge \textbf{会放大原有错误}（模型判不出自己错在哪，反而确信）→ 必须留失败案例
\end{enumerate}

\subsection{四、里程碑}\label{ux56dbux91ccux7a0bux7891}

\begin{itemize}
\tightlist
\item[$\square$]
  M1 手写最小 ReAct 循环（无框架），跑通单步任务
\item[$\square$]
  M2 加入第二个工具 + 多步任务集（≥30 个任务）
\item[$\square$]
  M3 \textbf{轨迹记录规范}：每步的 thought/action/observation 落盘（可复现）
\item[$\square$]
  M4 基线成功率 + \textbf{失败五分类}（每类给 2 个真实案例）
\item[$\square$]
  M5 \textbf{self-judge 消融}：三组对照，给出成功率差与误差棒
\item[$\square$]
  M6 \textbf{反例收集}：找出 ≥3 个''self-judge 导致更差''的案例并解释机制
\item[$\square$]
  M7 一页报告 + 知识卡
\end{itemize}

\subsection{五、交付物}\label{ux4e94ux4ea4ux4ed8ux7269}

{\def\LTcaptype{none} % do not increment counter
\begin{longtable}[]{@{}lll@{}}
\toprule\noalign{}
交付物 & 位置 & 验收 \\
\midrule\noalign{}
\endhead
\bottomrule\noalign{}
\endlastfoot
ReAct 循环 & \texttt{手写/} & 可跑通多步任务 \\
self-judge 协议 & \texttt{手写/}（R10） & 有明确的判定时机与阈值设计说明 \\
轨迹日志 & \texttt{证据/trajectories/} & 每任务一份，含步数与结果 \\
消融数据 & \texttt{证据/} & 三组成功率 + 误差棒 \\
报告 & \texttt{报告.md} & 必须回答''何时 self-judge 有害'' \\
\end{longtable}
}

\subsection{六、讲课锚点}\label{ux516dux8bb2ux8bfeux951aux70b9}

\begin{itemize}
\tightlist
\item
  \textbf{讲义}：\texttt{M8-智能体.md}（W13 展开）
\item
  \textbf{对标}：\textbf{李宏毅《Context Engineering》专讲}（Agent 时代的关键技术）· Berkeley LLM Agents MOOC（ReAct 原始论文那讲）· HF Agents Course（三框架对照 + 观测性/评测 bonus unit）· \textbf{rasbt ch5（self-refinement）} · \texttt{nanochat/execution.py}+\texttt{tests/test\_execution.py}（\textbf{沙箱执行的最小实现}，可读接口不可抄实现）
\item
  \textbf{搭桥}：M5 的 KV cache（长上下文的显存代价）→ 为什么长程 Agent 必须做\textbf{上下文压缩}
\end{itemize}

\subsection{七、代码级提问示例}\label{ux4e03ux4ee3ux7801ux7ea7ux63d0ux95eeux793aux4f8b}

\begin{enumerate}
\def\labelenumi{\arabic{enumi}.}
\tightlist
\item
  长程任务失败率随步数指数上升。\textbf{分解}：是每步错误率问题，还是错误会累积传播？
\item
  self-judge 什么时候\textbf{一定}无效？（提示：当模型无法区分''看起来对''和''真的对''时）
\item
  上下文压缩（截断/摘要）会引入什么新失败模式？
\item
  为什么''工具返回值要结构化''能显著提升成功率？（从 token 分布角度解释）
\item
  \textbf{【下注】} 预测 self-judge 在简单任务 vs 困难任务上的效果差异，先写再测。
\end{enumerate}

\subsection{八、风险与提醒}\label{ux516bux98ceux9669ux4e0eux63d0ux9192}

\begin{itemize}
\tightlist
\item
  🚨 \textbf{不要用框架起步}：手写一遍 ReAct 再去看 LangGraph/smolagents，否则学不到 harness 的本质。
\item
  ⚠️ self-judge 会显著增加 token 消耗与时间 → 算力预算要写进实验卡。
\item
  ⚠️ 轨迹日志可能很大 → 只保留结构化摘要（thought 截断 + 结果），避免污染仓库体积。
\end{itemize}

\end{document}
